\section{Check if a Linked List is a Palindrome}
\label{sec:ll-palindrome}

\problemstatement{Given the head of a singly linked list, determine whether
the sequence of values reads the same forwards and backwards, using $O(1)$
extra space.}

\begin{patternbox}
Without random access, comparing front to back means \textbf{finding the
middle (slow/fast), reversing the second half, then walking both halves in
step}. Equality of the two halves proves a palindrome. This combines two
earlier techniques into one.
\end{patternbox}

\subsection*{Middle, reverse, compare}

Use the tortoise--hare to reach the middle. Reverse the second half in place
(the iterative reversal). Now walk one pointer from the head and another
from the reversed second half, comparing values; if all match, the list is a
palindrome. Optionally restore the list by reversing the second half back.
This achieves $O(1)$ space, versus copying values into an array ($O(n)$).

\begin{ideabox}[Reverse Half to Compare Without Extra Memory]
A palindrome check needs to read the list backward, which a singly linked
list cannot do -- so reverse the back half physically. Comparing the
untouched front half against the reversed back half is exactly the
front-to-back comparison, in constant space.
\end{ideabox}

\subsection*{Algorithm}

\begin{enumerate}
  \item Find the middle with slow/fast.
  \item Reverse the second half.
  \item Compare the first half and the reversed second half node by node.
  \item Return whether all comparisons matched.
\end{enumerate}

\subsection*{Dry run}

\dryrun{List $1\to2\to2\to1$.}

\begin{itemize}
  \item Middle after node $2$ (second); reverse back half $\to1\to2$.
  \item Compare $1{=}1$, $2{=}2$ -- palindrome.
\end{itemize}

\subsection*{C++ implementation}

\begin{lstlisting}
#include <bits/stdc++.h>
using namespace std;

struct ListNode {
    int val; ListNode* next;
    ListNode(int v) : val(v), next(nullptr) {}
};

ListNode* build(const vector<int>& a) {
    ListNode dummy(0), *t = &dummy;
    for (int x : a) { t->next = new ListNode(x); t = t->next; }
    return dummy.next;
}

ListNode* reverse(ListNode* head) {
    ListNode* prev = nullptr;
    while (head) { ListNode* nx = head->next; head->next = prev; prev = head; head = nx; }
    return prev;
}

bool isPalindrome(ListNode* head) {
    ListNode* slow = head;
    ListNode* fast = head;
    while (fast && fast->next) { slow = slow->next; fast = fast->next->next; }
    ListNode* second = reverse(slow);             // reverse back half

    ListNode* p1 = head;
    ListNode* p2 = second;
    while (p2) {
        if (p1->val != p2->val) return false;
        p1 = p1->next; p2 = p2->next;
    }
    return true;
}

int main() {
    cout << isPalindrome(build({1, 2, 2, 1})) << "\n";   // 1
    cout << isPalindrome(build({1, 2, 3})) << "\n";      // 0
    return 0;
}
\end{lstlisting}

\begin{complexitybox}
\textbf{Time Complexity.} $O(n)$ -- find middle, reverse, compare, each
linear. \textbf{Space Complexity.} $O(1)$.
\end{complexitybox}

\begin{takeawaybox}
Palindrome checking composes slow/fast (middle) with in-place reversal
(back half) to compare front and back in constant space. It is the classic
demonstration that combining two basic list techniques beats the $O(n)$-
memory array copy.
\end{takeawaybox}
